\subsection{Complete m-adic rings and inverse limits}

We have seen in no. 6 that, if A is a commutative ring and m an ideal of A such that A is Hausdorff and complete with the m-adic topology, then the topological ring A is canonically identified with the inverse limit of the discrete rings \(A_{i} = A/m^{i+1} (i \in \mathbf{N})\) with respect to the canonical mappings

\[
h _ {i j} \colon \mathrm{A} / \mathfrak {m} ^ {j + 1} \rightarrow \mathrm{A} / \mathfrak {m} ^ {i + 1} \quad (i \leqslant j);
\]

note that \(h_{ij}\) is surjective and that, if \(n_{ij}\) is its kernel, then

\[
\mathfrak {n} _ {i j} = \mathfrak {m} ^ {i + 1} / \mathfrak {m} ^ {j + 1} = (\mathfrak {m} / \mathfrak {m} ^ {j + 1}) ^ {i + 1} = (\mathfrak {n} _ {0 j}) ^ {i + 1};
\]

in particular \((\mathfrak{n}_{0j})^{j + 1} = 0\) . Conversely:

PROPOSITION 14. Let \((\mathbf{A}_i, h_{ij})\) be an inverse system of discrete commutative rings, whose indexing set is \(\mathbf{N}\) and let \((\mathbf{M}_i, u_{ij})\) be an inverse system of modules over the inverse system of rings \((\mathbf{A}_i, h_{ij})\) . Let \(\mathfrak{n}_j\) denote the kernel of \(h_{0j} \colon \mathbf{A}_j \to \mathbf{A}_0\) and set \(\mathbf{A} = \varprojlim \mathbf{A}_i\) , \(\mathbf{M} = \lim \mathbf{M}_i\) . Suppose that

\begin{enumerate}
\def\labelenumi{(\alph{enumi})}
\item
  for all \(i \in \mathbf{N}\) , \(h_{i1}\) is the identity mapping on \(A_i\) and, for \(i \leqslant j\) , \(h_{i1}\) and \(u_{i1}\) are surjective;
\item
  for \(i \leqslant j\) , the kernels of \(h_{ij}\) and \(u_{ij}\) are \(\mathfrak{n}_j^{i+1}\) and \(\mathfrak{n}_j^{i+1}\mathrm{M}_j\) respectively.
\end{enumerate}

Then:

\begin{enumerate}
\def\labelenumi{(\roman{enumi})}
\item
  A is a complete Hausdorff topological ring, M is a complete Hausdorff topological A-module and the canonical mappings \(h_{i}: A \to A_{t}, u_{i}: M \to M_{i}\) are surjective.
\item
  If \(M_{0}\) is a finitely generated \(A_{0}\) -module, M is a finitely generated A-module;

to be precise, every finite subset S of M such that \(u_{0}(S)\) generates \(M_{0}\) is a system of generators of M.
\end{enumerate}

The assertions in (i) follow from General Topology, Chapter II, § 3, no. 5, Corollary to Proposition 10 and Corollary 1 to Theorem 1.

For all \(i \in \mathbf{N}\) , let us write \(\mathfrak{m}_{i+1} = \operatorname{Ker}(h_i)\) , \(\mathbf{N}_{i+1} = \operatorname{Ker}(u_i)\) ; then

\[
\mathfrak {m} _ {i + 1} = \varprojlim_ {k \geqslant 0} h _ {i, i + k} ^ {- 1} (0) = \varprojlim_ {k} \mathfrak {n} _ {i + 1} ^ {i + k}
\]

and \(\mathbf{N}_{i + 1} = \lim_{k\to \infty}\mathfrak{n}_{i + k}^{i + 1}\mathbf{M}_{i + k}\) ; as \(h_{i + k}\) and \(u_{i + k}\) are surjective,

\[
h _ {i + k} \left(\mathfrak {m} _ {i + 1}\right) = \mathfrak {n} _ {i + k} ^ {i + 1}, \quad u _ {i + k} \left(\mathrm{N} _ {i + 1}\right) = \mathfrak {n} _ {i + k} ^ {i + 1} \mathrm{M} _ {i + k}.\tag{16}
\]

Let us show that \(\mathfrak{m}_i\mathrm{N}_j\subset \mathrm{N}_{i + j}\) for \(i\geqslant 1\) and \(j\geqslant 1\) , which amounts to proving that \(u_{i + j - 1}(\mathfrak{m}_i\mathrm{N}_j) = 0\) ; now

\[
u _ {i + j - 1} (m _ {i} \mathrm{N} _ {j}) = h _ {i + j - 1} (m _ {i}) u _ {i + j - 1} (\mathrm{N} _ {j})
\]

is equal to \(\mathfrak{n}_{i+j-1}^{i}(\mathfrak{n}_{i+j-1}^{j}\mathrm{M}_{i+j-1}) = 0\) , as, for all \(k \geqslant 0\) , \(\mathfrak{n}_k^{k+1}\) , which is the kernel of \(h_{kk}\) , is equal to 0. We see similarly that \(\mathfrak{m}_i\mathfrak{m}_j \subset \mathfrak{m}_{i+j}\) . If for \(i \leqslant 0\) we set \(\mathfrak{m}_i = \mathrm{A}\) and \(\mathrm{N}_i = \mathrm{M}\) , \((\mathfrak{m}_i)_{i \in \mathbf{Z}}\) is a filtration of \(\mathrm{A}\) and \((\mathrm{N}_i)_{i \in \mathbf{Z}}\) a filtration of \(\mathrm{M}\) compatible with the filtration on \(\mathrm{A}\) ; the topologies on \(\mathrm{A}\) and \(\mathrm{M}\) are obviously those defined by these filtrations. This being so, let \(\mathfrak{a}\) be an ideal of \(\mathrm{A}\) such that \(h_1(a) = \mathfrak{n}_1\) and \(\mathrm{M}'\) the submodule of \(\mathrm{M}\) generated by \(\mathrm{S}\) ; we are going to prove that

\[
\mathrm{N} _ {i} = \mathfrak {a} ^ {i} \mathrm{M} ^ {\prime} + \mathrm{N} _ {i + 1} \quad \text { for } \quad i \geqslant 0.\tag{17}
\]

Let us write \(a_i = h_i(a), M'_i = u_i(M')\) ; it is sufficient to show that

\[
u _ {i} \left(\mathrm{N} _ {i}\right) = a _ {i} ^ {t} \mathrm{M} _ {i} ^ {\prime}.
\]

This is true if \(i = 0\) , for \(\mathbf{N}_0 = \mathbf{M}\) and \(\mathbf{M}_0' = \mathbf{M}_0\) by hypothesis. If \(i \geqslant 1\) , then \(u_i(\mathbf{N}_i) = \mathfrak{n}_i^i\mathbf{M}_i\) by (16). As \(h_{1i}\) is surjective and \(h_{0i} = h_{01} \circ h_{1i}\) , \(h_{1i}\) maps the kernel \(\mathfrak{n}_i\) of \(h_{0i}\) onto the kernel \(\mathfrak{n}_1\) of \(h_{01}\) and \(\mathfrak{n}_t = \overline{h}_{1i}^{-1}(\mathfrak{n}_1)\) ; then

\[
h _ {1 i} (a _ {i}) = h _ {1} (a) = n _ {1} = h _ {1 i} (n _ {i})
\]

and, as the kernel of \(h_{1i}\) is \(n_i^2\) , \(n_i \subset a_i + n_i^2\) and \(a_i \subset n_i\) , whence \(n_i = a_i + n_i^2\) . Moreover \(u_{0i}(M_i') = u_0(M') = M_0 = u_{0i}(M_i)\) and, as \(\text{Ker}(u_{0i}) = n_iM_i\) , \(M_t = M_t' + n_iM_t\) ; whence

\[
n _ {i} ^ {t} \mathbf {M} _ {i} = (a _ {i} + n _ {i} ^ {2}) ^ {t} (\mathbf {M} _ {i} ^ {\prime} + n _ {i} \mathbf {M} _ {i}).
\]

Now, \(\mathfrak{a}_i^k\mathfrak{n}_i^{i + 1 - k}\subset \mathfrak{n}_i^{i + 1} = 0\) for \(0\leqslant k\leqslant i\) ; then it certainly follows that \(u_{i}(\mathbf{N}_{i}) = \mathfrak{n}_{i}\mathbf{M}_{i} = \mathfrak{a}_{i}^{i}\mathbf{M}_{t}^{\prime}\) , which proves (17).

Moreover \(m_1 = h_1(n_1)\) , whence \(a \subset m_1\) and therefore \(a^t \subset m_1^t \subset m_i\) , whence

\(N_{i} \subset m_{i}M' + N_{i+1}\) ; on the other hand obviously \(m_{i}M \subset N_{i}\) and hence \(N_{i} = m_{i}M' + N_{i+1}\) for all \(i \geqslant 0\) ; then it follows from Corollary 3 to Proposition 12 of no. 9 that \(M' = M\) , which completes the proof.

COROLLARY 1. With the notation and hypotheses of Proposition 14 suppose further that \(\mathbf{M}_0\) is a finitely generated \(\mathbf{A}_0\) -module and that the ideal \(\mathfrak{n}_1\) of \(\mathbf{A}_1\) is finitely generated. Let \(\mathfrak{m}_1\) be the kernel of \(h_0\) ; the topologies on \(\mathbf{A}\) and \(\mathbf{M}\) are then the \(\mathfrak{m}_1\) -adic topologies on this ring and this module respectively; to be precise, for all \(i \geqslant 0\) , the kernels of \(h_i\) and \(u_i\) are \(\mathfrak{m}_1^{i+1}\) and \(\mathfrak{m}_1^{i+1}\mathbf{M}\) respectively; further \(\mathfrak{m}_1/\mathfrak{m}_1^2\) is a finitely generated A-module.

We preserve the notation of the proof of Proposition 14; the hypotheses here allow us to assume that the ideal a is finitely generated. Let \(i \geqslant 0\) be any integer; for all \(j \geqslant 0\) , by (17), \(\mathbf{N}_{i+j} = \mathfrak{a}^j (\mathfrak{a}^i\mathbf{M}) + \mathbf{N}_{i+j+1} \subset \mathfrak{m}_j(\mathfrak{a}^i\mathbf{M}) + \mathbf{N}_{i+j+1}\) ; conversely, \(\mathfrak{m}_j(\mathfrak{a}^i\mathbf{M}) \subset \mathfrak{m}_j\mathfrak{m}_i\mathbf{M} \subset \mathfrak{m}_{i+j}\mathbf{M} \subset \mathbf{N}_{i+j}\) , whence

\[
\mathbf {N} _ {i + j} = \mathfrak {m} _ {j} (\mathfrak {a} ^ {i} \mathbf {M}) + \mathbf {N} _ {i + j + 1}.
\]

As a and M are finitely generated A-modules, so is \(a^{i}M\) . Applying Corollary 3 to Proposition 12 of no. 9 to the module \(N_{i}\) with the filtration \((\mathrm{N}_{ij})_{j\in\mathbf{Z}}\) defined by \(N_{ij}=N_{i}\) if j\textless0, \(N_{ij}=N_{i+j}\) if \(j\geqslant0\) , we obtain \(N_{i}=a^{i}M\) , whence \(N_{i}\subset m_{1}^{i}M\) . But also \(m_{1}^{i}M\subset m_{i}M\subset N_{i}\) , whence \(N_{i}=m_{1}^{i}M\) . Applying this to the case where \(M_{i}=A_{i}\) , \(u_{ij}=h_{ij}\) , we obtain \(m_{i}=m_{1}^{i}\) . Moreover, \(m_{1}=a+m_{1}^{2}\) by (17), which proves the last assertion of the corollary.

COROLLARY 2. Under the hypotheses of Corollary 1, for A to be Noetherian, it is necessary and sufficient that \(A_{0}\) be so.

The condition is necessary since \(A_{0}\) is isomorphic to a quotient of A; it is sufficient by virtue of no. 10, Corollary 5 to Theorem 2.

